\documentclass[11pt]{amsart}

\usepackage[T1]{fontenc}
\usepackage{lmodern}
\usepackage{microtype}
\usepackage{amsmath,amssymb,booktabs}
\usepackage[hidelinks]{hyperref}

\newcommand{\perev}{\operatorname{per}^{\mathrm{ev}}}
\newcommand{\Oev}{\Omega^{\mathrm{ev}}}
\newcommand{\Sn}{\mathfrak S}
\newcommand{\An}{\mathfrak A}

\newtheorem{theorem}{Theorem}

\title[An order-six Brualdi--Liu counterexample]
{An order-six counterexample to the Brualdi--Liu even-permanent conjecture}

\date{}

\subjclass[2020]{15A15}
\keywords{even permanent, even permutation polytope, doubly stochastic matrix}

\begin{document}

\begin{abstract}
Brualdi and Liu conjectured that, for $n\ge 3$, every matrix $X$ in the convex
hull $\Oev_n$ of the even permutation matrices satisfies
$\perev(X)\ge n!/(2n^n)$. Wanless gave counterexamples for $n=4$ and
$n=5$ and left open the possibility that the inequality might hold at larger
orders. We settle the next order by giving an explicit $M\in\Oev_6$ with
$\perev(M)=661/100000<5/648$. The inequality also fails on the relative
interior of $\Oev_6$.
\end{abstract}

\maketitle

For $\sigma\in\Sn_n$, let $P_\sigma$ be the permutation matrix with a $1$ in
position $(i,\sigma(i))$, and set
\[
  \perev(X)=\sum_{\sigma\in\An_n}\prod_{i=1}^n x_{i,\sigma(i)},
  \qquad
  \Oev_n=\operatorname{conv}\{P_\sigma:\sigma\in\An_n\}.
\]
Brualdi and Liu conjectured that, for $n\ge 3$ and $X\in\Oev_n$,
\begin{equation}\label{eq:BL}
  \perev(X)\ge \frac{n!}{2n^n},
\end{equation}
with equality only at the matrix whose entries are all $1/n$
\cite[Conjecture~2]{BrualdiLiu}. Wanless disproved \eqref{eq:BL} for
$n=4$ and $n=5$ and left open whether it might hold for larger $n$
\cite{Wanless}. We settle the next order.

\begin{theorem}\label{thm:main}
The Brualdi--Liu inequality fails for $n=6$. More precisely, there exists
$M\in\Oev_6$ such that
\[
  \perev(M)=\frac{661}{100000}<\frac{5}{648}
  =\frac{6!}{2\cdot6^6}.
\]
Moreover, \eqref{eq:BL} fails at points in the relative interior of
$\Oev_6$.
\end{theorem}

\begin{proof}
Let
\[
A=
\begin{pmatrix}
3&2&5&0&0&0\\
0&3&2&3&0&2\\
0&0&3&4&3&0\\
2&3&0&3&2&0\\
5&0&0&0&0&5\\
0&2&0&0&5&3
\end{pmatrix},
\qquad
M=\frac1{10}A.
\]
Define
\[
\begin{aligned}
\pi_1&=(2\,4)(5\,6),
&\pi_2&=(1\,2\,3\,4)(5\,6),\\
\pi_3&=(1\,3\,5),
&\pi_4&=(1\,3\,4\,5)(2\,6).
\end{aligned}
\]
These permutations are even, and direct inspection gives
\[
  A=3P_{\pi_1}+2P_{\pi_2}+3P_{\pi_3}+2P_{\pi_4}.
\]
Hence $M\in\Oev_6$.

Writing $\operatorname{inv}(\sigma)$ for the inversion number, a direct
expansion over the support of $A$ gives exactly six nonzero even terms:
\[
\begin{array}{c@{\qquad}c@{\qquad}r}
\toprule
(\sigma(1),\ldots,\sigma(6))
  & \operatorname{inv}(\sigma)
  & \displaystyle\prod_{i=1}^6 a_{i,\sigma(i)}\\
\midrule
(1,3,4,5,6,2) & 4  & 480\\
(1,4,3,2,6,5) & 4  & 2025\\
(2,3,4,1,6,5) & 4  & 800\\
(2,3,4,5,1,6) & 4  & 480\\
(3,2,5,4,1,6) & 6  & 2025\\
(3,6,4,5,1,2) & 10 & 800\\
\bottomrule
\end{array}
\]
Thus
\[
  \perev(M)
  =\frac{480+2025+800+480+2025+800}{10^6}
  =\frac{661}{100000}.
\]
The strict inequality follows from
\[
  661\cdot648=428328<500000=5\cdot100000.
\]

Let $J_6$ be the all-ones matrix. Since
\[
  \frac{J_6}{6}
  =\frac1{360}\sum_{\sigma\in\An_6}P_\sigma,
\]
for every $t\in(0,1)$ the matrix
\[
  M_t=(1-t)M+t\frac{J_6}{6}
\]
has a positive coefficient on every vertex of $\Oev_6$, and hence belongs
to its relative interior. Since $\perev$ is continuous and
$\perev(M)<5/648$, the same strict inequality holds for all sufficiently
small $t>0$.
\end{proof}

\begin{thebibliography}{9}

\bibitem{BrualdiLiu}
Richard A.~Brualdi and Bolian~Liu,
\emph{The polytope of even doubly stochastic matrices},
J. Combin. Theory Ser. A \textbf{57} (1991), no.~2, 243--253,
\href{https://doi.org/10.1016/0097-3165(91)90048-L}
{doi:10.1016/0097-3165(91)90048-L}.

\bibitem{Wanless}
Ian M.~Wanless,
\emph{On the Brualdi--Liu conjecture for the even permanent},
Electron. J. Linear Algebra \textbf{17} (2008), 284--286,
\href{https://doi.org/10.13001/1081-3810.1263}
{doi:10.13001/1081-3810.1263}.

\end{thebibliography}

\end{document}